% Adapted from the official SimulacraBench technical report template.
\documentclass{article}
\PassOptionsToPackage{numbers,sort&compress}{natbib}
\usepackage[final]{neurips_2026}
\makeatletter
\providecommand{\@trackname}{40th Conference on Neural Information Processing Systems (NeurIPS 2026). SimulacraBench Competition.}
\makeatother
\usepackage[utf8]{inputenc}
\usepackage[T1]{fontenc}
\usepackage{booktabs,amsmath,amssymb,url,hyperref}
\hypersetup{hidelinks}
\title{Estimating Survey Skip Probabilities and\\Combining Predictors by Question}
\author{Appalanaidu Bobbili}
\begin{document}
\maketitle
\begin{center}\small\begin{tabular}{l}\toprule
\textbf{Team ID:} Survey Routing \quad (Test submission 1057678)\\\bottomrule
\end{tabular}\end{center}

\begin{abstract}
Survey questionnaires use skip rules to determine who receives a follow-up.
We estimate the probability of being asked from visible survey records, allowing
exceptions to the prescribed rule. When the preceding answer is also hidden, we
average over its possible values, including its own skip outcome. For each
question, we use held-out visible respondents to select a mixture of this
predictor and a categorical naive Bayes predictor. In a synthetic experiment
with deliberately imperfect rules, estimating skip probabilities increased
mean log score from $-1.579$ to $-0.910$ nats per cell against the same answer
model with hard rules. Under exact rules both methods scored $-0.510$ to three
decimals. On a separate synthetic draw, the mixture improved mean normalized
skill from 0.2844 to 0.2930 versus naive Bayes alone. These comparisons used
a uniform question prior. The submitted mixture achieved a displayed overall
skill score of 0.35 on the official SimulacraBench Test evaluation.
\end{abstract}

\section{Method summary}
\subsection{Method type}
\begin{tabular}{@{}ll@{}}
$\boxtimes$ Bayesian / hierarchical model & $\boxtimes$ Imputation\\
$\boxtimes$ Supervised prediction & $\boxtimes$ Ensemble of multiple components\\
\end{tabular}\\[3pt]
An optional language model supplies a prior over the substantive answers.

\subsection{Pipeline}
Consider a follow-up prescribed after a positive response. Two probabilities
matter: how often the respondent receives the follow-up, and which answer they
give if asked. We estimate both. The first predictor uses the organizer's
hierarchical answer model~\cite{starter} and replaces its deterministic skip
rule with rates estimated from visible records. The second predicts every
outcome, including the skip, directly from visible attributes. A held-out
comparison determines how much weight each predictor receives for each item.

\subsection{Inputs and feature construction}
Both predictors use categorical \texttt{GIVEN} attributes and the complete
responses available in the instrument. We read the answer levels, their order,
and the skip-rule parent from the schema. A skip is a scored category distinct
from refusals and other recorded responses. The respondent identifier is
excluded. A skipped value on a visible attribute has its own feature level.
Only the question wording, options and public survey description
enter the optional language-model prompt.

\newpage
\subsection{Model or simulator}
For parent level $k$, let $n_k$ count visible parent--child pairs and $a_k$
count how often the child was asked. Group levels by the prescribed state
$b\in\{0,1\}$ (closed or open), with pooled counts $N_b$ and $A_b$. We use
\[
 r_b=\frac{A_b+8b+0.5}{N_b+9},\qquad r_k=\frac{a_k+8r_b}{n_k+8}.
\]
A rare level borrows evidence from levels with the same prescribed state.
If the parent is hidden, $P(\mathrm{asked}\mid x)=\sum_k P(k\mid x)r_k$ uses
the full predicted parent distribution, including its own skip when applicable.
The conditional answer predictor shrinks counts from up to two visible
attributes toward successively coarser strata.

\subsection{Fitting, training, or prompting}
For each instrument we mask every target answer of a random fifth of the
complete visible respondents (seed 2026). Both predictors use the remaining
rows. For each question we choose a mixing weight by the mean log probability
of the observed answers on these masked rows, then refit on all visible
respondents. Appendix~\ref{config} specifies the weight grid and the inner
validation used for naive Bayes.

\subsection{Prediction and aggregation}
For a gated item, the first predictor multiplies its probability of each
substantive answer by $P(\mathrm{asked}\mid x)$, then appends the skip
probability. The final vector mixes this distribution with the second
predictor's full distribution. Each masked cell contributes one vector in row
and schema order. Respondents receive equal weight within an item.

\subsection{Post-processing and calibration}
The outer weight search applies the grader's floor $\varepsilon=0.001$ in
the logarithmic score~\cite{scores}. We return normalized vectors before
flooring. The grader applies the floor once.

\section{Models, data, and compute}
\begin{center}\small
\begin{tabular}{@{}ll@{}}\toprule
Component & Role and provenance\\\midrule
Qwen2.5-1.5B-Instruct~\cite{qwen} & Optional item prior, offline weights, no fine-tuning\\
SituatedEvals external-model example~\cite{starter} & Conditional answer counts and prompt construction\\
Categorical naive Bayes & Shrunk attribute likelihoods and a marginal fallback\\
\bottomrule\end{tabular}
\end{center}
\paragraph{Data.} Local evaluations use the official
schema-shaped generator and a two-question toy generator included with the
code. In the phase-two rehearsal, the synthetic training and development
respondents are visible, and every target of the synthetic test respondents
is masked.\\[3pt]
\begin{tabular}{@{}llll@{}}\toprule External dataset & Source & Size & Use\\\midrule
None & --- & --- & ---\\\bottomrule\end{tabular}
\paragraph{Compute.} The CPU rehearsals in Appendix~\ref{results-local} use
Python 3.10.11, NumPy 2.2.6, pandas 2.3.3, PyYAML 6.0.3 and PyTorch
2.14.1+cpu. They invoke the uniform-prior branch. We also tested Qwen on an
RTX 5060 Laptop GPU with PyTorch 2.7.1, CUDA 12.8 and Transformers 4.53.3.
The tokenizer and language model both require public
checkpoint revision \texttt{989aa7980e4cf806f80c7fef2b1adb7bc71aa306}.

\newpage
\section{Results}
\subsection{Test-phase leaderboard results}
The official Test leaderboard lists \texttt{1057678} for Survey Routing
(\texttt{submission\_routed.zip}). We report its four displayed scores at the
leaderboard's precision~\cite{leaderboard}.\\[4pt]
\begin{tabular}{@{}lrrrr@{}}\toprule Split & Grand Prize & UNICEF & World Bank & UNHCR\\\midrule
Test & 0.35 & 0.33 & 0.21 & 0.52\\\bottomrule\end{tabular}
\subsection{Breakdown by population, survey, or question type}
The highest skill was on UNHCR (0.52) and the lowest was on World Bank (0.21).
All three instrument scores exceeded the smoothed marginal reference of zero.
These aggregates do not identify which respondent groups or questions account
for the differences.

\section{Reproducibility}
The upstream starter revision is \texttt{5c69e17}. The team implementation
is in \texttt{submission\_routed/} and carries the MIT license. From the
repository root, these commands rebuild and validate the competition archive.
\begin{small}\begin{verbatim}
python tools/pack_candidate.py submission_routed submission_routed.zip
python tools/check_submission_zip.py submission_routed.zip
\end{verbatim}\end{small}
The SHA-256 digest of the submission archive is\\
{\small\texttt{0bd3220a562815c8568c6943dbe9ad639274f94389c1f14fa7fc31bca997f84e}}.
The source commit that reproduces this archive is\\
{\small\texttt{50e60fd6af92cebe699e361f736270db8f38f646}} in the team repository:\\
\url{https://github.com/ab-naidu/neurips-2026-simulacrabench-survey-routing}.
Appendix~\ref{config} identifies the model and configuration.

\section{Limitations}
The public generator enforces skip rules exactly. The UNHCR schema documents
exceptions and partial overlap between expected branches~\cite{starter}.
The stress experiment measures the cost of those kinds of deviations at
invented rates. We do not know the frequency of departures on the Test respondents.
A single validation split can give unstable weights for rare outcomes.
Individual predictions do not establish representativeness or subgroup accuracy.

\begin{thebibliography}{9}
\bibitem{starter} SituatedEvals. \emph{SimulacraBench public starter kit},
including the external-model baseline and UNHCR schema, 2026.
\url{https://github.com/SituatedEvals/public/tree/5c69e17}.
\bibitem{scores} T. Gneiting and A. E. Raftery. Strictly proper scoring rules,
prediction, and estimation. \emph{JASA}, 102(477), 359--378, 2007.
\bibitem{qwen} Qwen Team. \emph{Qwen2.5-1.5B-Instruct}, model card and weights.
\url{https://huggingface.co/Qwen/Qwen2.5-1.5B-Instruct}.
\bibitem{leaderboard} SimulacraBench. \emph{Codabench Test leaderboard},
submission 1057678, November 17, 2026.
\url{https://www.codabench.org/api/phases/29773/get_leaderboard/?format=json}.
\end{thebibliography}

\clearpage
\appendix
\section{Full configuration}\label{config}
\paragraph{Answer distributions.} Starting with the question prior, update
observed answer counts with two pseudo-observations from the preceding level.
Apply this update to the whole instrument, the highest-ranked attribute, then
the pair of highest-ranked attributes. Rank by mutual information with the
answer among visible, asked rows. Fewer than eight usable observations gives
zero mutual information. The gate calculation uses two fixed shrinkage
strengths of eight. Use the marginal ask rate $(a+0.5)/(n+1)$ for an
unmodeled parent or a missing visible parent value. Masked target parents
use their predicted distribution.
\paragraph{Direct predictor.} Use one half-count per outcome for the marginal.
Shrink each attribute distribution conditional on the outcome toward its
observed attribute frequencies with strength eight. Add the log marginal to
the sum of attribute log likelihoods scaled by one of
$\{0,.05,.1,.2,.35,.5,.75,1\}$. Choose the scale on every fifth visible row
in the expert's fitting set by mean log score, clipping probabilities below
$10^{-12}$ for numerical stability. Use scale
zero with fewer than 100 visible rows, no attributes, or a gain below .001
nats per cell.
\paragraph{Mixture.} For each question choose the gate predictor's weight from
$\{1,.75,.5,.25,0\}$. With $K$ outcomes and probability $q_y$ assigned to the
observed outcome, maximize the mean of $\log(.001+(1-.001K)q_y)$. Scores
within $10^{-12}$ of the best favor the larger gate weight. With fewer than
100 complete visible rows, use the gate predictor alone.
\paragraph{Optional question prior.} The hosted builder reads
\texttt{Qwen/Qwen2.5-1.5B-Instruct} from \texttt{models.txt}. The model
uses the offline cache and the checkpoint's chat template. For each question
it receives the following single user message, without demonstrations.
\begin{small}\begin{verbatim}
You are estimating how a survey respondent answered.

Survey: {first 600 characters of the public survey description}

Question: {schema question text}

Options:
A. {first option}
B. {second option}
{continue in schema order, up to Z}

Reply with the single letter of the most likely answer.
\end{verbatim}\end{small}
For each letter use the larger first-token logit of the bare and space-prefixed
letter. Softmax over letters. Additional options receive the mean option
probability before renormalization. If the chat template fails, append a
newline and \texttt{Answer:} to the message. If the model is unavailable,
the prior is uniform.

\newpage
\section{Additional results on synthetic data}\label{results-local}
Table~\ref{tab:practice} compares candidates on the same synthetic respondents
for each seed. Seed 0 was used during development. Seed 17 was generated after
the routing grid was fixed. We repeated the UNICEF evaluation after correcting
the encoding of skips in a visible attribute. Both seeds use the final-phase split
and scorer. Following the October 26 scoring update, instrument skill is
$1-L/L_{\mathrm{marginal}}$, where each $L$ is mean negative log probability
on the same masked cells and the reference adds a half-count per category.
The mean gives each instrument equal weight.
\begin{center}\small
\begin{tabular}{@{}llrrrr@{}}
\toprule
Seed & Model & UNICEF & World Bank & UNHCR & Mean\\\midrule
0 & Direct predictor & .3801 & .2621 & .2145 & .2856\\
0 & Organizer hard gates & .3693 & .2395 & .1482 & .2523\\
0 & Estimated gates & .3693 & .2395 & .1482 & .2523\\
0 & Per-question mixture & .3882 & .2723 & .2152 & .2919\\\midrule
17 & Direct predictor & .3710 & .2712 & .2110 & .2844\\
17 & Per-question mixture & .3838 & .2837 & .2115 & .2930\\
\bottomrule
\end{tabular}
\refstepcounter{table}\label{tab:practice}\\[3pt]
Table \thetable. Local normalized skill with a uniform question prior.
\end{center}
Across both seeds the largest gains over the direct predictor occurred on
UNICEF and World Bank. On UNHCR the gains were below .001 skill. The sum of
the three CPU grader times for the mixture was about 84 s on seed 0 and 80 s
on seed 17. These are single measurements including local grader overhead.

In a separate offline load check, Qwen produced valid priors for all 301
questions on an RTX 5060 Laptop GPU, with nonuniform priors on 296 questions.
With 120 visible and one masked respondent per schema, PyTorch reserved at
most 3,474 MiB of GPU memory. On the full seed-0 rehearsal the Qwen-enabled
mixture scored .3882 for UNICEF, .2723 for World Bank and .2152 for UNHCR.
All three match the uniform-prior mixture to four decimals. The grader
covered 442,425 masked cells in 504 s across three separate processes,
including a model load in each. The official scores do not identify which
prior ran on the hosted H100 or its runtime and peak memory.


The paired stress test isolates the skip model. It uses three fixed seeds
(41--43), two questions, 4,000 complete respondents and 2,000 masked
respondents per seed. The first question has a visible parent. The second
has a hidden parent that can itself have been skipped. The conditional answer
model and uniform prior are identical in both arms. Higher log score is better.
\begin{center}\small
\begin{tabular}{@{}lrr@{}}\toprule
Skip-rule regime & Hard rule (nats/cell) & Estimated rates (nats/cell)\\\midrule
Exact & $-0.510160$ & $-0.510161$\\
Noisy & $-1.579152$ & $-0.909969$\\\bottomrule
\end{tabular}\end{center}
In the noisy regime the probability of asking the first question is .85 or .22
depending on eligibility. For the second it is .81, .12 or .31 given a positive
answer, a negative answer or a skipped parent. Both parent types improved on
all three seeds. Run \texttt{python tests/test\_gate\_stress.py} to reproduce
this comparison.
\end{document}
